\section{Coverage and Access Diagnostics}
\label{app:inventory}

\begin{table}[t]
\centering\small
\begin{tabular}{lrr}
\toprule
Skill family & Test tasks & SFT examples\\
\midrule
\inputtable{tables/inventory.tex}
\bottomrule
\end{tabular}
\caption{Unified benchmark and training inventory. Family coverage is complete
but sizes and task distributions are not balanced or proportional. Counts do
not establish coverage of every principal operation within a Skill.}
\label{tab:family-inventory}
\end{table}

The five families with three benchmark tasks only are emergent projects,
motivation diagnosis, emotional literacy, self-explanation, and spaced
practice. Lesson building contributes 154 source scenarios and three
constructed tasks. The underlying task objects and order are unchanged between
the frozen export and the present scoring archive, despite JSON formatting
differences in their file hashes.

\begin{table*}[t]
\centering\small
\begin{tabular}{lrrrr}
\toprule
Skill family & GLM Flash & DS V4.1 Flash & DS V4 Pro & Doubao 2.1 Pro\\
\midrule
\inputtable{tables/skill_routes.tex}
\bottomrule
\end{tabular}
\caption{Completed expected-package loads / available tasks. Unknown captures
are retained in the available-task denominator. Absence of a competing
educational package means these counts measure access, not full-library
retrieval accuracy.}
\label{tab:family-access}
\end{table*}

\subsection{Complete Model--Skill Score Breakdown}

\begin{table*}[t]
\centering\small
\setlength{\tabcolsep}{4pt}
\begin{tabular}{lrrrrrr}
\toprule
Skill family & $n_s$ & GLM-5.3 & GLM Flash & DS Flash & DS Pro & Doubao\\
\midrule
\inputtable{tables/skill14_delta_all.tex}
\bottomrule
\end{tabular}
\caption{With--No score differences across all 14 Skills. Every family
retains its full task denominator, including execution failures. These are
the numerical values behind Figure~\ref{fig:skill-breakdown}; additional
decimals expose rounding, not statistical certainty.}
\label{tab:family-main}
\end{table*}

Tables~\ref{tab:skill14-detail-a} and~\ref{tab:skill14-detail-b} retain the
No/With means behind every difference in Table~\ref{tab:family-main}.
Each mean uses all $n_s$ tasks in that family, including zero-valued execution
failures. The same fixed task membership is used for all five models.
Table~\ref{tab:skill14-failures} reports the corresponding execution failures;
a completed answer that receives a low rubric score is not an execution
failure. Per-task answer bindings, both judgment components, common-valid
denominators, paired-direction counts, and weighted contributions to the
305-task mean are retained in the offline analysis release.

\begin{table*}[t]
\centering\small
\setlength{\tabcolsep}{3pt}
\begin{tabular}{lrrrrrrrrrr}
\toprule
& & \multicolumn{3}{c}{GLM-5.3} & \multicolumn{3}{c}{GLM Flash} & \multicolumn{3}{c}{DS V4.1 Flash}\\
\cmidrule(lr){3-5}\cmidrule(lr){6-8}\cmidrule(lr){9-11}
Skill family & $n_s$ & No & With & $\Delta$ & No & With & $\Delta$ & No & With & $\Delta$\\
\midrule
\inputtable{tables/skill14_scores_a.tex}
\bottomrule
\end{tabular}
\caption{Complete fourteen-Skill breakdown for the three result inventories
shown: reference-relative No/With means on the 0--100 scale and differences
in score points. Deltas are calculated from unrounded means. Scores are not
scientific-accuracy percentages.}
\label{tab:skill14-detail-a}
\end{table*}

\begin{table*}[t]
\centering\small
\setlength{\tabcolsep}{4pt}
\begin{tabular}{lrrrrrrr}
\toprule
& & \multicolumn{3}{c}{DS V4 Pro} & \multicolumn{3}{c}{Doubao Seed 2.1 Pro}\\
\cmidrule(lr){3-5}\cmidrule(lr){6-8}
Skill family & $n_s$ & No & With & $\Delta$ & No & With & $\Delta$\\
\midrule
\inputtable{tables/skill14_scores_b.tex}
\bottomrule
\end{tabular}
\caption{Complete fourteen-Skill breakdown for Pro and Doubao, with the same
denominators and scoring definition as Table~\ref{tab:skill14-detail-a}.
Pro's No value is 50 in every family because it is the reference-self
condition; it is not an independent judgment of Pro answer quality.}
\label{tab:skill14-detail-b}
\end{table*}

\begin{table*}[t]
\centering\small
\setlength{\tabcolsep}{4pt}
\begin{tabular}{lrrrrrr}
\toprule
Skill family & $n_s$ & GLM-5.3 & GLM Flash & DS Flash & DS Pro & Doubao\\
\midrule
\inputtable{tables/skill14_failures.tex}
\bottomrule
\end{tabular}
\caption{Execution failures in No/With order for every model--Skill group.
These counts expose completion-driven changes without removing failures
from the reported scores. A zero count does not mean the answers are
correct, or that the designated package was used.}
\label{tab:skill14-failures}
\end{table*}

\subsection{Evidence Levels and Case Selection}

\begin{table}[t]
\centering\scriptsize
\setlength{\tabcolsep}{3pt}
\begin{tabular}{lrrrr}
\toprule
Model & Loaded & Rate (\%) & $\Delta_L$ & $\Delta_N$\\
\midrule
\inputtable{tables/access_summary_v4.tex}
\bottomrule
\end{tabular}
\caption{Designated-package loads among 305 With tasks. $\Delta_L$ and
$\Delta_N$ are With--No changes on common-valid pairs with a load and with
no recorded access. DeepSeek Flash and Doubao each have one unknown
capture, retained in rate denominators.}
\label{tab:routing}
\label{tab:access-effects}
\end{table}

\begin{table}[t]
\centering\scriptsize
\setlength{\tabcolsep}{3pt}
\begin{tabular}{lrrrr}
\toprule
& \multicolumn{2}{c}{Loaded} & \multicolumn{2}{c}{No recorded access}\\
\cmidrule(lr){2-3}\cmidrule(lr){4-5}
Model & $N$ & $\Delta$ & $N$ & $\Delta$\\
\midrule
\inputtable{tables/access_effects.tex}
\bottomrule
\end{tabular}
\caption{Common-valid subgroup denominators and With--No score differences
for Table~\ref{tab:routing}. Groups are selected by the agent's recorded
access behavior; their differences do not identify a causal loading effect.}
\label{tab:access-subgroups}
\end{table}

Raw-trace access requires a completed designated Skill call with returned
content in the root answer's execution tree. Summary-level access requires a
retained completed-call record with the designated name but permits less
inspection of subsequent resource actions. The two unknown captures concern
DeepSeek Flash's \texttt{lesson-builder\_\_cn01\_05} and Doubao's
\texttt{lesson-builder\_\_03}. They are not counted as proven non-access.
One No Skill GLM Flash trace loads a built-in OpenCode customization package;
this is not an exposed educational Skill and is not counted as such.

The process-analysis bundle includes 32 case exports with saved No/With
answers, judgment components, and available tool evidence. Cases illustrate
positive, negative, tied, and incomplete outcomes; they do not estimate a
population error rate. The native records distinguish main-package loading
from supporting-file reads. In DeepSeek Flash, the 60 read errors comprise 41
pedagogy, 16 checkpoint, two quality-gate, and one course-passport reference
failures. Their frequency is specific to this mounted resource inventory.

\subsection{Candidate and Reference Component Changes}
\label{app:components}

\begin{table*}[t]
\centering\small
\begin{tabular}{llrrrr}
\toprule
Model & Access group & $N$ & $\Delta J_{\mathrm{cand}}$ & $\Delta J_{\mathrm{ref}}$ & $\Delta R$\\
\midrule
\inputtable{tables/score_components.tex}
\bottomrule
\end{tabular}
\caption{With--No changes on common valid pairs grouped by With access.
All three deltas are multiplied by 100; $\Delta R$ averages per-task ratios
rather than applying the ratio to the two mean components. Pro's reference-self
No condition has no judged component scores, shown as unavailable rather than
imputed.}
\label{tab:components}
\end{table*}

The same reference text may receive different component scores in separate
candidate pairs. Positive ratio changes can therefore accompany decreases in
the judged reference component. Retaining these quantities prevents
normalization from being mistaken for an independent measurement of absolute
improvement. Neither component analysis nor anonymous presentation proves
immunity to position, length, self-family, or reference-quality biases.
